\section{Phản ứng xà phòng hóa và sản xuất xà phòng}

\subsection{Phản ứng của chất béo với kiềm}

Khi đun nóng chất béo với dung dịch kiềm như \ce{NaOH}, chất béo bị thủy phân tạo glycerol và muối của acid béo. Phản ứng này được gọi là phản ứng xà phòng hóa.

\[
\ce{(RCOO)3C3H5 + 3NaOH ->[t^\circ] 3RCOONa + C3H5(OH)3}.
\]

Muối sodium hoặc potassium của một số acid béo là thành phần chính của xà phòng.

Acid béo tự do cũng phản ứng với kiềm:
\[
\ce{RCOOH + NaOH -> RCOONa + H2O}.
\]

\begin{ghinho}{Ghi nhớ}
- Chất béo phản ứng với kiềm tạo glycerol và muối của acid béo.
- Muối sodium của acid béo thường là thành phần chính của xà phòng rắn.
- Phản ứng xà phòng hóa là phản ứng thủy phân chất béo trong môi trường kiềm.
\end{ghinho}

\subsection{Bài toán sản xuất xà phòng}

Trong bài toán sản xuất xà phòng, cần liên hệ lượng chất béo hoặc acid béo ban đầu với lượng muối thu được theo phương trình hóa học. Nếu nguyên liệu có tạp chất hoặc quá trình có hiệu suất, phải tính đến các dữ kiện đó.

\begin{vidu}
Cho \SI{89.0}{\gram} tristearin \ce{(C17H35COO)3C3H5} phản ứng vừa đủ với dung dịch \ce{NaOH}, đun nóng. Tính khối lượng sodium stearate \ce{C17H35COONa} thu được. Cho $M_{\ce{(C17H35COO)3C3H5}}=890$ g/mol và $M_{\ce{C17H35COONa}}=306$ g/mol.
\end{vidu}
\begin{loigiai}
\[
n_{\ce{(C17H35COO)3C3H5}}=\frac{89.0}{890}=\SI{0.100}{\mole}.
\]

\bangmol{7}
{\ce{(C17H35COO)3C3H5} & + & \ce{3NaOH} & \rightarrow & \ce{3C17H35COONa} & + & \ce{C3H5(OH)3}}
{0.100 & & & & 0.300 & &}

\[
m_{\ce{C17H35COONa}}=0.300\times306=\SI{91.8}{\gram}.
\]

\textbf{Vậy,} khối lượng sodium stearate thu được là \SI{91.8}{\gram}.
\end{loigiai}

\subsection{Chỉ số acid và chỉ số ester hóa}

Trong một số bài toán nâng cao về chất béo, đề có thể cho:
\begin{itemize}
\item \textbf{Chỉ số acid}: số miligam \ce{KOH} cần dùng để trung hòa acid béo tự do có trong \SI{1}{\gram} chất béo.
\item \textbf{Chỉ số ester hóa}: số miligam \ce{KOH} cần dùng để xà phòng hóa lượng triglyceride có trong \SI{1}{\gram} chất béo.
\end{itemize}
Khi đề đã cho hai chỉ số này, có thể đổi khối lượng \ce{KOH} tiêu thụ thành số mol, sau đó dùng tỉ lệ phản ứng để xác định lượng acid béo tự do và triglyceride. Đây là dữ kiện mở rộng, không phải công thức cần học thuộc riêng.
